%%%%%%%%%%%%%%%%%%%%%%%%%%%%%%%%%%%%%%%%%%%%%%%%%%%%%%%%%%%%%
%% Variables %%%%%%%%%%%%%%%%%%%%%%%%%%%%%%%%%%%%%%%%%%%%%%%%
%%%%%%%%%%%%%%%%%%%%%%%%%%%%%%%%%%%%%%%%%%%%%%%%%%%%%%%%%%%%%

\newcommand{\primarycolor}[1]{\def\theprimarycolor{#1}}
\newcommand{\secondarycolor}[1]{\def\thesecondarycolor{#1}}
\newcommand{\counterbackground}[1]{\def\thecounterbackground{#1}}

%%%%%%%%%%%%%%%%%%%%%%%%%%%%%%%%%%%%%%%%%%%%%%%%%%%%%%%%%%%%%
%% Colors %%%%%%%%%%%%%%%%%%%%%%%%%%%%%%%%%%%%%%%%%%%%%%%%%%%
%%%%%%%%%%%%%%%%%%%%%%%%%%%%%%%%%%%%%%%%%%%%%%%%%%%%%%%%%%%%%

\definecolor{darkcerulean}{rgb}{0.063661, 0.257392, 0.477463}
\definecolor{scooter}{rgb}{0.161162, 0.775760, 0.885416}
\definecolor{fireenginered}{rgb}{0.757398, 0.151664, 0.177088}
\definecolor{pomegranate}{rgb}{0.927557, 0.260794, 0.138377}
\definecolor{mineshaft}{rgb}{0.199974, 0.200015, 0.199971}
\definecolor{tingray}{rgb}{0.501904, 0.501994, 0.501898}


\primarycolor{darkcerulean}
\secondarycolor{scooter}

%%%%%%%%%%%%%%%%%%%%%%%%%%%%%%%%%%%%%%%%%%%%%%%%%%%%%%%%%%%%%
%% Beamer Layout %%%%%%%%%%%%%%%%%%%%%%%%%%%%%%%%%%%%%%%%%%%%
%%%%%%%%%%%%%%%%%%%%%%%%%%%%%%%%%%%%%%%%%%%%%%%%%%%%%%%%%%%%%
\useoutertheme[subsection=false]{miniframes}
\useinnertheme{default}
\usefonttheme[onlymath]{serif}
\setbeamercolor*{normal text}{fg=black,bg=white}
\setbeamercolor*{section in head/foot}{fg=white, bg=\theprimarycolor}
\setbeamercolor*{frametitle}{fg=\theprimarycolor, bg=white}
\setbeamercolor*{titlelike}{fg=\theprimarycolor,bg=white}
\setbeamercolor*{structure}{fg=\theprimarycolor}
 
\setbeamertemplate{logo}{
    \vspace{-1em}
    \begin{tikzpicture}
        \node {\thecounterbackground};
        \node {{\tiny\color{white}\insertframenumber}};
    \end{tikzpicture}
}
\setbeamertemplate{footline}{}
\beamertemplatenavigationsymbolsempty

\setbeamertemplate{itemize items}{
    \begin{tikzpicture}
        \fill (0,0) circle (2pt);
        \draw (0,0) circle (3pt);
    \end{tikzpicture}
}

\setbeamertemplate{enumerate items}[default]


%****************************************************************************
% The following is a hack to allow the transparent package to be loaded at the
% same time as the tikz package: 
% https://tex.stackexchange.com/questions/253401/tikz-and-ctable-incompatibility-gives-error-when-printing/253417#253417


\makeatletter
\begingroup\expandafter\expandafter\expandafter\endgroup
\expandafter\ifx\csname pgfutil@addpdfresource@extgs\endcsname\relax
\else
  \AtBeginDocument{%
    % \pgf@sys@addpdfresource@extgs@plain{%
    \pgfutil@addpdfresource@extgs{%
      \TRP@list
    }%
  }%  
  \let\TRP@addresource\relax
\fi
\makeatother
%****************************************************************************




%***************************************************************************************************
% The following magic makes sure that the pnote command produces notes that are displayed in pdfpc:

% create a new file handle
\newwrite\pdfpcnotesfile

% open file on \begin{document}
\AtBeginDocument{%
	\immediate\openout\pdfpcnotesfile\jobname.pdfpc\relax
	\immediate\write\pdfpcnotesfile{[notes]}
}
% define a # http://tex.stackexchange.com/a/37757/10327
\begingroup
	\catcode`\#=12
	\gdef\hashchar{#}%
\endgroup


\def\lastframenumber{0}

% define command \pnote{} that works like note but
% additionally writes notes to file in pdfpc readable format
\newcommand{\pnote}[1]{%
	% keep normal notes working
	\note{#1}%

	% if frame changed - write a new header
	\ifdim\theframenumber pt>\lastframenumber pt
		\let\lastframenumber\theframenumber
		\begingroup
			\let\#\hashchar
			\immediate\write\pdfpcnotesfile{\#\#\# \theframenumber}%
		\endgroup
	\fi

	% write note to file
	\immediate\write\pdfpcnotesfile{\unexpanded{#1}}%
}
% close file on \begin{document}
\AtEndDocument{%
	\immediate\closeout\pdfpcnotesfile
}

%            END OF pdfpc MAGIC
%***************************************************************************************************

%%%%%%%%%%%%%%%%%%%%%%%%%%%%%%%%%%%%%%%%%%%%%%%%%%%%%%%%%%%%%
%%% Boxes
%%%%%%%%%%%%%%%%%%%%%%%%%%%%%%%%%%%%%%%%%%%%%%%%%%%%%%%%%%%%%


% #1 – box name
% #2 – left color
% #3 – right color
% #4 – Title (may contain #1)
\newcommand{\newlargebox}[4]{
    \newtcolorbox{#1}[1][]{
        beamer,
        width=\textwidth,
        titlerule=3mm,
        colframe=white,
        coltitle={#2},
        bottom=4pt, 
        top=-4pt,
        left=6pt,   
        right=6pt,   
        toptitle=4pt,
        bottomtitle=-2pt,
        fonttitle=\bfseries\large,
        adjusted title={#4},
        outer arc=.5mm,
        arc=.5mm,
        no shadow,
        fuzzy shadow={1mm}{-1mm}{-1.2mm}{.7mm}{black!20},
        interior titled code={
          \path [left color = {#2}, right color = {#3}]
          (title.south west) + (8pt, 0) rectangle ++(\textwidth-8pt, 0.02);
        }
      }
}


\newtcolorbox{emptybox}[1][]{
    beamer,
    width=\textwidth,
    % enlarge left by=-3pt,
    titlerule=0mm,
    colframe=white,
    coltitle=black,
    bottom=8pt, 
    top=-10pt,
    left=6pt,  
    notitle, 
    adjusted title={},
    outer arc=.5mm,
    arc=.5mm,
    no shadow,
    fuzzy shadow={1mm}{-1mm}{-1.2mm}{.7mm}{black!20},
    interior titled code={}
}

\newtcolorbox{plainbox}[1][]{
    beamer,
    width=\textwidth,
    % enlarge left by=-3pt,
    titlerule=0mm,
    colframe=white,
    coltitle=black,
    bottom=8pt,
    toptitle=4pt, 
    top=2pt,
    left=6pt,  
    notitle, 
    adjusted title={#1},
    outer arc=.5mm,
    arc=.5mm,
    no shadow,
    fuzzy shadow={1mm}{-1mm}{-1.2mm}{.7mm}{black!20},
    interior titled code={}
}

\newlargebox{defaultbox}{\theprimarycolor}{\thesecondarycolor}{#1} 
\newlargebox{darkbox}{mineshaft}{tingray}{#1} 
\newlargebox{graybox}{tingray}{white}{#1} 

\newlargebox{theorembox}{\theprimarycolor}{\thesecondarycolor}{Theorem {\ifthenelse{\equal{#1}{}}{}{(#1)}}} 
\newlargebox{definitionbox}{\theprimarycolor}{\thesecondarycolor}{Definition {\ifthenelse{\equal{#1}{}}{}{(#1)}}} 
\newlargebox{proofbox}{mineshaft}{tingray}{Proof{\ifthenelse{\equal{#1}{}}{}{ #1}}.} 

\newcommand{\smallbox}[1]{\begin{emptybox}\textbf{#1}\end{emptybox}}
\newcommand{\defaultsmallbox}[1]{\smallbox{\color{\theprimarycolor}#1}}
\newcommand{\graysmallbox}[1]{\smallbox{\color{tingray}#1}}

